\section{Measurement Suite}
\label{sec:measurement}
\textbf{Support fidelity} includes source-relative appearance and geometric displacement, with registration failures reported. \textbf{Material behavior} separates activity, active area and direction where observable. \textbf{Continuity} targets replay and resets while retaining legitimate periodic flow. \textbf{Scene integrity} addresses missing or invented content beyond the measured support. Human ratings assess the corresponding observable attributes separately. Figure~\ref{fig:routing} shows which role each measure reads.
% Which pixels each measure reads. Definitions: docs/METRIC_SPEC_v1.1.md; realism family: Section IV-A.
% Type floor: \scriptsize (7 pt in IEEEtran journal 10 pt); nothing smaller. Fits \columnwidth (8.89 cm).
\begin{figure}[t]
\centering
\begin{tikzpicture}[font=\scriptsize,
  role/.style={draw=#1,fill=#1!18,line width=.7pt,rounded corners=1.5pt,text width=1.42cm,
    minimum height=6mm,align=center,inner sep=1.5pt},
  meas/.style={draw=snfink!60,fill=white,rounded corners=1.5pt,text width=3.2cm,align=left,inner sep=2pt},
  real/.style={draw=snfink!60,fill=snfpaper,rounded corners=1.5pt,text width=2.92cm,align=left,inner sep=2pt},
  wire/.style={->,>={Stealth[length=1.3mm]},line width=.55pt,#1}]
\node[role=roleSupport]  (sup) at (0,0)     {\textbf{Support}\\[-1pt]should stay fixed};
\node[role=roleMaterial] (mat) at (0,-2.2)  {\textbf{Material}\\[-1pt]should keep moving};
\node[role=roleSubject]  (sub) at (0,-3.3)  {\textbf{Subject}\\[-1pt]v2.1, planned};
\node[role=roleOverlay]  (ovl) at (0,-4.1)  {\textbf{Overlay}};
\node[role=roleIgnore]   (ign) at (0,-4.75) {\textbf{Uncertain}};

\node[meas,anchor=west] (m1) at (0.98,0) {\gsupport\ \textbf{Support fidelity}\\
  \stfrozen\ fBD\down\ displacement\\ \stfrozen\ NBF\down\ background flow\\ \stcand\ source $\Delta E_{00}$\down};
\node[meas,anchor=west,fill=snfpaper] (m3) at (0.98,-1.1) {\stfrozen\ DLR (context):\\ support\,/\,material};
\node[meas,anchor=west] (m2) at (0.98,-2.2) {\gmotion\ \textbf{Material motion}\\
  \stfrozen\ MCFF-E/L flow\\ \stfrozen\ FP persistence\\ \stfrozen\ DAR (context)};
\node[meas,anchor=west] (m4) at (0.98,-3.3) {\gintegrity\ \stopen\ Scene integrity};
\node[anchor=west,text=snfink!75,text width=3.2cm,inner sep=2pt] (m5) at (0.98,-4.42)
  {\gnone\ enter no measure;\\ reported as coverage};

\node[real,anchor=west] (r1) at (4.55,0)     {\gsupport\ \stdiag\ R1 support integrity};
\node[real,anchor=west] (r3) at (4.55,-1.55) {\gdirection\ \stcand\ R3 direction};
\node[real,anchor=west] (r4) at (4.55,-2.2)  {\greplay\ \stcand\ R4 replay};
\node[real,anchor=west] (r5) at (4.55,-2.85) {\gplaus\ \stopen\ R5 plausibility};
\node[real,anchor=west] (r2) at (4.55,-3.5)  {\gmaterial\ \stopen\ R2 presence};

\draw[wire=roleSupport]  (sup.east) -- (m1.west);
\draw[wire=roleSupport]  (sup.east) -- ($(m3.west)+(0,0.06)$);
\draw[wire=roleMaterial] (mat.east) -- (m2.west);
\draw[wire=roleMaterial] (mat.east) -- ($(m3.west)+(0,-0.06)$);
\draw[wire=roleSubject]  (sub.east) -- (m4.west);
\draw[wire=snfgrey]      (ovl.east) -- ($(m5.west)+(0,0.1)$);
\draw[wire=snfgrey]      (ign.east) -- ($(m5.west)+(0,-0.1)$);
\draw[wire=roleSupport]  (m1.east) -- (r1.west);
\foreach \r in {r3,r4,r5,r2} \draw[wire=roleMaterial] (m2.east) -- (\r.west);
\draw[wire=roleSubject]  (m4.east) -- (r2.west);
\node[font=\footnotesize\bfseries,anchor=south] at ($(sup.north)+(0,0.12)$) {Role};
\node[font=\footnotesize\bfseries,anchor=south] at ($(m1.north)+(0,0.05)$) {Axes and context};
\node[font=\footnotesize\bfseries,anchor=south] at ($(r1.north)+(0,0.05)$) {Realism family};
\end{tikzpicture}
\caption{\textbf{Illustrative role routing.} Lines show intended inputs, not complete computational dependencies: DLR compares roles and motion compensation can use support information. Scene integrity and subject handling are planned; R2 is material presence. Overlay and uncertain pixels are excluded from the current partitioned measures and reported as coverage. Realism
measures are reported beside the axes and never summed with them. A downward arrow marks a directional support statistic, not an overall quality score; material activity and the contextual ratios require joint interpretation. \stfrozen\ frozen definition, construct validation pending; \stcand\ development candidate;
\stdiag\ diagnostic only; \stopen\ planned. Interactive version: \texttt{interactive/measurement\_map.html}.}
\label{fig:routing}
\end{figure}


These are candidate axes, not an admitted headline score. A sharp static frame can preserve detail without motion; noise can increase activity; natural surf can exhibit periodicity. Frame-based displacement (fBD) remains a diagnostic, not a sufficient surrogate for source fidelity. A historical mean-colour field has a documented float32 accumulation defect and retains its original measurement identity.

\subsection{Realism Family}
Realism is a separate family of measures, reported beside support fidelity, material behavior and continuity, and never summed with them. Each measure names the single construct it targets. A measure that orders a statistic in one direction can reward an exploit: activity rewards incoherent motion, and closeness to the source rewards a pasted still image. Each realism statistic is therefore read relative to an explicit reference: real fixed-camera footage of the same medium where the construct allows, otherwise encoded controls.
\begin{itemize}
\item \emph{Support integrity.} Successive-frame change in the support after drift and global-illumination compensation, read against real footage and against pass-through, pasted, grain-matched and feathered composites. It is a diagnostic: a valid fixed background can be still and a careful composite can look natural, so it does not by itself establish that a support was pasted or redrawn.
\item \emph{Material presence and location.} Whether the requested material, rather than another region or a subject, carries the motion. It is checked against source-conditioned material and subject annotations; for example, a moving boat over frozen water is not persistent water flow.
\item \emph{Direction.} The flow component along an image-plane direction that is annotated independently of the scene text, with unknown, ambiguous and multi-direction options, defined only where the direction is visible, and tested against time-reversed and ping-pong controls of matched magnitude.
\item \emph{Progression versus replay.} Multi-frame recurrence with seam and reset evidence, gated by active area. Detection and false-positive rates are reported separately, with natural periodic surf, fire and foliage as negatives.
\item \emph{Motion plausibility.} Flow, point-track and dynamic-texture statistics inside the dynamic role, compared with medium-level distributions of real footage with source-family uncertainty, and with established comparators on the same eligible inputs: TRAJAN per video~\cite{allen2025trajan}, FVMD at set level only~\cite{liu2024fvmd}, and VBench-2.0 physics~\cite{Zheng2025VBench20AV} with any adaptation recorded.
\end{itemize}
A real continuation of a source frame is one plausible realization, not the correct trajectory, so no measure compares pixels with it. No two-dimensional statistic establishes physical correctness, and none of these measures is claimed validated before Section~\ref{sec:validation}. A scientifically plausible cinemagraph can be a valid output; a composite counts as a defect only where the task or human raters say so.

\textbf{Eligibility and abstention.} A measure abstains when its region is empty or registration to the source fails, and the number of eligible videos is reported beside every value. Abstentions are never imputed, and registration failures are counted rather than dropped.
